\section{Newtonian three-dimensional bubble}
\label{sec:newtonian}

\status{Uniform mesh is \textbf{closed at both resolutions} --- $25\times25\times50$
and the reference $50\times50\times100$, 1501 steps each, in both methods. Case~2
has not been run. The adaptive half now has a working dynamic remeshing loop
(Section~\ref{sec:amr}) but no production run.}

\subsection{Case definition}

A bubble of radius $R = 0.25$ centred at $(0.5,0.5,0.5)$ in a tank
$[0,1]\times[0,1]\times[0,2]$, no-slip on all walls, zero initial velocity,
pressure and polymeric tensor, integrated to $t = 3$. Two parameter sets,
Table~\ref{tab:params}, taken from the reference work and originally due to
Adelsberger et al.

\begin{table}[htbp]
\centering
\caption{Physical and dimensionless parameters.}
\label{tab:params}
\small
\pgfplotstabletypeset[
  col sep=space, string type,
  every head row/.style={before row=\toprule, after row=\midrule},
  every last row/.style={after row=\bottomrule},
]{\dat{caso-parametros.dat}}
\end{table}

\subsection{Uniform mesh}

\subsubsection{The methods converge toward each other}

This is the central result of the section, and it is a statement about both
methods rather than about either one: refining the mesh brings them together.

\begin{table}[htbp]
\centering
\caption{Difference between the two methods at $t=3$, at both resolutions.}
\label{tab:convergence}
\small
\begin{adjustbox}{max width=\linewidth}
\pgfplotstabletypeset[
  col sep=space, string type,
  columns/grandeza/.style={column name={Quantity}, column type={l}},
  columns/n25/.style={column name={$25\times25\times50$}, column type={r}},
  columns/n50/.style={column name={$50\times50\times100$}, column type={r}},
  columns/fator/.style={column name={Factor}, column type={r}},
  every head row/.style={before row=\toprule, after row=\midrule},
  every last row/.style={after row=\bottomrule},
]{\dat{c1-convergencia.dat}}
\end{adjustbox}
\end{table}

What makes this more than a coincidence is the \emph{direction} each method
moves. The VOF centre of mass falls from $1.4903$ to $1.4766$; the
front-tracking one rises from $1.4615$ to $1.4678$. They approach from opposite
sides toward roughly $1.47$. A single method converging says only that it is
consistent; two independent representations converging from opposite sides
toward the same value is evidence that the value is the right one.

\begin{table}[htbp]
\centering
\caption{Front-tracking against VOF at $t=3$, both resolutions. Shape
quantities from the extracted surfaces by the same routine; mass from the field
integral for VOF and from the surface volume for front-tracking.}
\label{tab:resolution}
\small
\begin{adjustbox}{max width=\linewidth}
\pgfplotstabletypeset[
  col sep=space, string type,
  columns/grandeza/.style={column name={Quantity at $t=3$}, column type={l}},
  columns/vof25/.style={column name={VOF $N{=}25$}, column type={r}},
  columns/ft25/.style={column name={FT $N{=}25$}, column type={r}},
  columns/vof50/.style={column name={VOF $N{=}50$}, column type={r}},
  columns/ft50/.style={column name={FT $N{=}50$}, column type={r}},
  every head row/.style={before row=\toprule, after row=\midrule},
  every last row/.style={after row=\bottomrule},
]{\dat{c1-resolucao.dat}}
\end{adjustbox}
\end{table}

Both methods reproduce the \textbf{ellipsoidal regime} the reference reports for
case~1 --- the height-to-equivalent-diameter ratio is below one in both, that
is, flattened along the direction of rise --- and both reproduce the shape of
the published rise-velocity curve: the velocity grows, peaks near $t = 0.9$, and
decreases toward a steady value.

\begin{figure}[htbp]
\centering
\begin{tikzpicture}
\begin{axis}[width=0.47\linewidth, xlabel={$t$}, ylabel={$z_c$},
             legend pos=north west, grid=both, no markers, legend cell align=left]
  \addplot table[x=t, y=zc] {\dat{bolha3d-c1-N25-vof.dat}}; \addlegendentry{VOF $N{=}25$}
  \addplot table[x=t, y=zc] {\dat{bolha3d-c1-N25-ft.dat}};  \addlegendentry{FT $N{=}25$}
  \addplot table[x=t, y=zc] {\dat{bolha3d-c1-N50-vof.dat}}; \addlegendentry{VOF $N{=}50$}
  \addplot table[x=t, y=zc] {\dat{bolha3d-c1-N50-ft.dat}};  \addlegendentry{FT $N{=}50$}
\end{axis}
\end{tikzpicture}\hfill
\begin{tikzpicture}
\begin{axis}[width=0.47\linewidth, xlabel={$t$}, ylabel={$\Psi$},
             legend pos=south west, grid=both, no markers, legend cell align=left]
  \addplot table[x=t, y=Psi] {\dat{bolha3d-c1-N25-vof.dat}}; \addlegendentry{VOF $N{=}25$}
  \addplot table[x=t, y=Psi] {\dat{bolha3d-c1-N25-ft.dat}};  \addlegendentry{FT $N{=}25$}
  \addplot table[x=t, y=Psi] {\dat{bolha3d-c1-N50-vof.dat}}; \addlegendentry{VOF $N{=}50$}
  \addplot table[x=t, y=Psi] {\dat{bolha3d-c1-N50-ft.dat}};  \addlegendentry{FT $N{=}50$}
\end{axis}
\end{tikzpicture}
\caption{Centre of mass and sphericity at both resolutions.}
\label{fig:resolution}
\end{figure}

\subsubsection{Mass conservation goes the other way}

Mass is the one quantity that does \emph{not} improve for both methods, and the
two resolutions disagree about which method is better --- which is itself the
finding.

At $25\times25\times50$ the front-tracking drift is $-0.096\%$ against
$+0.396\%$ for VOF, a factor of four in favour of front-tracking. At
$50\times50\times100$ it reverses: $-0.851\%$ against $-0.170\%$, a factor of
five the other way.

Reading the time history explains it. Up to step 400 the two front-tracking runs
are nearly identical ($-7.3\times10^{-4}$ and $-6.8\times10^{-4}$ at step 200,
$-2.5\times10^{-3}$ and $-2.8\times10^{-3}$ at step 400). After that they
diverge: the coarse run peaks at $-2.6\times10^{-3}$ and \emph{recovers} to
$-9.6\times10^{-4}$, while the fine run accumulates monotonically. The VOF drift
likewise changes sign between the resolutions. Both coarse-mesh figures are
therefore partly \textbf{cancellation}, not conservation, and neither should be
quoted as a property of its method.

What survives is the fine-mesh statement: at the reference resolution VOF
conserves mass about five times better than front-tracking, and front-tracking's
drift \emph{grows} with refinement. The second half of that is the part worth
pursuing --- the triangulation ends with the same vertex and triangle counts at
both resolutions, and the surgery count differs by only 8.8\%, so neither
explains a factor of nine. \missing{the mechanism; candidate is the divergence
error of the velocity interpolated at the surface, which does not have to fall
with $h$ when the field it interpolates sharpens}

\paragraph{The instrument was cleared, not assumed.}
The VOF centre of mass is computed from the marching-cubes surface, which
encloses about 2.3\% less volume than the field at the coarse resolution; that
could bias the centroid. An independent probe computes the centroid directly
from the field and gives $1.490212$ against $1.4903$ from the surface ---
agreement to \num{1e-4}. The volume bias does not translate into a position
bias.

\paragraph{Topological soundness.}
121 output frames per method per resolution, \emph{all} surfaces closed (zero
orphan edges, $\chi = 2$). In the front-tracking runs the triangulation grew from
5120 to 5558 triangles, with 8824 and 9604 surgery operations and a
partition-of-unity residual of \num{2.1e-14}.

\subsection{Case 2: topological change}
\missing{$\rho_1/\rho_2 = 1000$, $\mu_1/\mu_2 = 100$, $\sigma = 1.96$. Lower
surface tension favours break-up, which is where an explicit surface with
surgery and a volume-fraction field should differ most. Not run.}

\subsection{Adaptive mesh}
\missing{blocked on Section~\ref{sec:amr}}
